\section{Preliminaries}\label{sec:preliminaries}

\paragraph{Parametric and contextual memory.}
We view a Transformer as coupling two retrieval mechanisms:
parameterized knowledge learned during training and contextual memory constructed in prompts
\citep{irie2022dual,zhong2025understandingtransformerperspectiveassociative}.
At layer $\ell$, the MLP processes a normalized hidden state $x\in\mathbb{R}^{d}$:
\begin{equation}
    m_\ell(x)=B_\ell\phi_\ell(x),\qquad
    \phi_\ell(x)=(A_\ell x)\odot\sigma(G_\ell x),
    \label{eq:mlp_assoc}
\end{equation}
where $A_\ell,G_\ell\in\mathbb{R}^{m\times d}$,
$B_\ell\in\mathbb{R}^{d\times m}$, and $\sigma$ is the
activation. Equivalently, $m_\ell(x)=\sum_{a=1}^{m}b_{\ell,a}
\phi_{\ell,a}(x)$, where $b_{\ell,a}$ is a column of $B_\ell$. So $\phi_\ell(x)$ means how an input addresses the
value vectors in $B_\ell$.

Attention retrieves values associated with the context. Let $q_{\ell,t}$,
$k_{\ell,s}$, and $v_{\ell,s}$ denote the query, key, and
output-projected value at token positions $t$ and $s$. Its output is
\begin{equation}
 a_{\ell,t}=\sum_{s\le t}\alpha_{\ell,ts}v_{\ell,s},\quad
 \alpha_{\ell,ts}=
 \frac{\exp(q_{\ell,t}^{\top}k_{\ell,s}/\sqrt{d_k})}
 {\sum_{u\le t}\exp(q_{\ell,t}^{\top}k_{\ell,u}/\sqrt{d_k})}.
 \label{eq:attn_assoc}
\end{equation}
The query-key similarity acts as a retrieval kernel,
while attention determines which earlier tokens can
contribute to later tokens.
The attention KV memories vary with the input and position, while
the MLP parameters persist as knowledge.

\paragraph{Writing and retrieving parameter updates.}
For a training example, let $\delta_{\ell,s}$ be the loss gradient
with respect to $m_\ell(x_{\ell,s})$ at position $s$.
A gradient-descent update of the MLP down-projection has the outer-product form
\begin{equation}
 \Delta B_\ell=-\eta\sum_s
 \delta_{\ell,s}\phi_\ell(x_{\ell,s})^{\top}.
 \label{eq:assoc_write}
\end{equation}
To determine which other inputs can read this update, we can calculate the MLP output delta on the gradient:
\begin{equation}
 \Delta m_\ell(x)
 =-\eta\sum_s\delta_{\ell,s}
 \underbrace{\phi_\ell(x_{\ell,s})^{\top}\phi_\ell(x)}
 _{K_\ell^{\mathrm{MLP}}(x_{\ell,s},x)}.
 \label{eq:assoc_transfer}
\end{equation}
It means that a training example changes the MLP output at another input in
proportion to their feature overlap.  In full-parameter training, evolving representations also change retrieval addresses.
Consequently, learning a fact does not guarantee its accessibility with a different query.

\paragraph{Compositional reasoning with knowledge updates.}
Consider two facts $e_2=r_1(e_1)$ and $e_3=r_2(e_2)$.
We denote the one-hop queries by $q_1=(e_1,r_1)$ and
$q_2=(e_2,r_2)$, and the two-hop query by
$q_{12}=(e_1,r_1,r_2)$. Both $q_2$ and $q_{12}$ request $e_3$,
but only $q_2$ explicitly contains $e_2$.
One possible compositional path forms a representation of $e_2$ at an
earlier layer, transports that information through
higher-layer attention, and activates
the association for $r_2(e_2)$. This mechanism is supported by studies of multi-hop reasoning
\citep{wang2024grokked,biran2024hopping}. 
If the latter association is read by an MLP in layer $j$, we can know whether the multi-hop state $x_j(q_{12})$ activates the same knowledge as the one-hop state $x_j(q_2)$, measured by $K_j^{\mathrm{MLP}}(x_j(q_{12}),x_j(q_2))$. A larger overlap indicates stronger activation of the features associated with the one-hop query, making the corresponding knowledge more accessible to the multi-hop query.

Updating the first fact, $r_1(e_1):e_2\mapsto e'_2$, requires the
two-hop query to produce $r_2(e'_2)$, assuming the downstream fact is
available. Updating the second fact, $r_2(e_2):e_3\mapsto e'_3$,
requires both $q_2$ and $q_{12}$ to produce $e'_3$.
The former changes the
intermediate representation supplied to downstream retrieval, whereas
the latter tests whether different queries reuse the same updated
association. Ideally, queries unrelated to the edited fact should retain their original answers. In practice, however, parameter updates may interfere with other knowledge, leading to unintended forgetting \citep{zhong2023mquake,cohen2023evaluatingrippleeffectsknowledge}.
